Put
\[
 D=\beta\kappa+2\beta+\kappa+1,
 \qquad e_n=\delta_{\mathrm{edge},n},
 \qquad c_n=\delta_{\mathrm{crit},n}.
\]
By the definitions of the two scales,
\[
 \frac{c_n}{e_n}
 =n^{1/(2\beta+\kappa+1)-1/(2D)}\longrightarrow\infty,
\]
because
\[
 2D-(2\beta+\kappa+1)=2\beta\kappa+2\beta+\kappa+1>0.
\]

Write \(T_n=\delta_n^{\kappa+1}h_n^\beta\) for the atom term in \(r_n\).  On every subsequence on which \(\delta_n/e_n\to\infty\), Lemma \(\texttt{lem:bandwidth-phases}\) gives
\[
 h_n\asymp(n\delta_n^\kappa)^{-1/(2\beta+1)}
\]
and hence
\[
 T_n\asymp
 n^{-\beta/(2\beta+1)}
 \delta_n^{D/(2\beta+1)}.
\]
Since \(c_n^D=n^{-1/2}\), this identity also yields
\[
 \frac{T_n}{n^{-1/2}}
 \asymp
 \left(\frac{\delta_n}{c_n}\right)^{D/(2\beta+1)}.
\]
On every subsequence on which \(\delta_n/e_n=O(1)\), the other conclusion of Lemma \(\texttt{lem:bandwidth-phases}\) gives \(h_n\asymp e_n\), and therefore
\[
 T_n\leq C e_n^{\beta+\kappa+1}
 =C n^{-(\beta+\kappa+1)/(2\beta+\kappa+1)}
 =o(n^{-1/2}),
\]
where the last relation follows from
\[
 \frac{\beta+\kappa+1}{2\beta+\kappa+1}-\frac12
 =\frac{\kappa+1}{2(2\beta+\kappa+1)}>0.
\]

Suppose first that \(\delta_n/c_n\to0\).  Every subsequence has a further subsequence on which \(\delta_n/e_n\) is bounded or tends to infinity.  The preceding two displays show on either such further subsequence that \(T_n=o(n^{-1/2})\).  The subsequence criterion for convergence therefore gives \(T_n=o(n^{-1/2})\) on the original sequence, so Definition \(\texttt{def:frontier}\) implies \(r_n\asymp n^{-1/2}\).

If \(\delta_n/c_n\to c\in(0,\infty)\), then \(c_n/e_n\to\infty\) implies \(\delta_n/e_n\to\infty\), and the displayed ratio for \(T_n\) is bounded above and away from zero.  Thus both summands of \(r_n\) have order \(n^{-1/2}\).  If \(\delta_n/c_n\to\infty\) and \(\delta_n\to0\), then again \(\delta_n/e_n\to\infty\), while the same ratio tends to infinity.  Consequently the atom term dominates and
\[
 r_n\asymp T_n
 \asymp\delta_n^{\kappa+1}(n\delta_n^\kappa)^{-\beta/(2\beta+1)}.
\]
If \(\delta_n\to\delta_0\in(0,\bar\delta]\), then \(\delta_n/e_n\to\infty\), so Lemma \(\texttt{lem:bandwidth-phases}\) yields \(h_n\asymp n^{-1/(2\beta+1)}\).  Because \(\delta_n^{\kappa+1}\) is then bounded above and away from zero and \(\beta/(2\beta+1)<1/2\),
\[
 r_n\asymp n^{-\beta/(2\beta+1)}.
\]
Theorem \(\texttt{thm:minimax-risk}\) and Theorem \(\texttt{thm:honest-length}\) transfer these comparisons of \(r_n\) to minimax absolute risk and minimax honest expected length, respectively.

It remains to verify the endpoint \(\delta_n=0\).  Proposition \(\texttt{prop:pushforward-setup}\) gives \(q_x(0)=0\) for every stratum.  Its pushforward formula at zero, together with \(d_0(A)=A\), also gives \(P(A=0\mid X=x)=0\); summing over the finitely many strata gives \(P(A=0)=0\).  The definition of the observed-data target therefore gives
\[
 \theta_0(P)=\mathbb E_P[Y\mathbf 1\{A>0\}]=\mathbb E_PY.
\]
In Definition \(\texttt{def:total-gram-estimator}\), \(\widehat b_x=0\) for every \(x\) almost surely, and \(\widehat\nu_n\) is the sample mean of the bounded outcomes on \(I_0\).  Thus the estimator center is exactly that bounded sample mean.  Finally, Definition \(\texttt{def:frontier}\) gives \(r_n=n^{-1/2}\) at \(\delta_n=0\), and the already established coverage and expected-length bounds in Theorem \(\texttt{thm:honest-length}\) show directly, without any extraneous equation reference, that the accompanying construction is honest bounded-mean inference of root-\(n\) order.  This proves every asserted regime.